\documentclass[12pt]{article}
\usepackage[utf8]{inputenc}
\usepackage[T1]{fontenc}
\usepackage{hyperref}
\title{Wom Language Resources for K-12 Curriculum Development}
\author{Compiled by Claude Sonnet 5.5 (claude-sonnet-5-5) for Andr\'as Kornai}
\date{}
\begin{document}
\maketitle

\section{Introduction}
Wom (also called Wam or Mie; ISO 639-3 wmo, Glottocode womp1235) is a Torricelli language of East Sepik Province, Papua New Guinea, with about 4,300 speakers. A 1975 survey places the language group roughly 22 km WNW of Maprik. Little published material was found.

\section{Primary Educational Resources}
No primers, textbooks or school materials were found.

\section{Linguistic and Reference Materials}
\begin{itemize}
\item Wikipedia article, with a phonology sketch: \url{https://en.wikipedia.org/wiki/Wom_language_(Papua_New_Guinea)}
\item Endangered Languages Project entry, which cites a 1981 sociolinguistic survey by Moeckel and Moeckel: \url{https://endangeredlanguages.com/lang/10565}
\end{itemize}
PARADISEC reportedly holds Wom material from Capell and Laycock, but no URL for it appeared in the search results, so it is not listed in the JSON. No dictionary or grammar was found.

\section{Religious and Community Resources}
Joshua Project reports that Bible translation into Wom has begun: \url{https://joshuaproject.net/languages/wmo}

\section{Digital and Audio Resources}
No online audio, video or software resources were found.

\section{Curriculum Development Context}
Wom is not known to be taught in schools, and the search results show no existing curriculum materials. A curriculum would need new work with the community, starting from the archival recordings and the survey cited above. Those sources have not been examined here.

\end{document}